% Teil II: der durchgerechnete Weg.

\subsection*{Gegeben und gesucht}
Gegeben: Der Läufer halbiert sein Trägheitsmoment, $I_2=\frac12 I_1$, ohne
äussere Drehmomente um die Achse.
Gesucht: Faktoren, mit denen sich (a) $L$, (b) $\omega$, (c) $E_\mathrm{rot}$
ändern.

\subsection*{Ansatz}
Drehimpulserhaltung um die Körperachse liefert $\omega_2$; damit die neue
Energie. Die Energieänderung aus der Arbeit der Armmuskeln begründen.

\subsection*{Rechnung}
\textbf{Drehimpuls $L$:} Ohne äussere Drehmomente ist $L$ streng erhalten:
\begin{equation}
  \frac{L_2}{L_1} = 1 \qquad (L\text{ ändert sich \emph{nicht}}).
\end{equation}

\textbf{Winkelgeschwindigkeit:} Aus $I_1\omega_1=I_2\omega_2$ und
$I_2=\frac12 I_1$:
\begin{equation}
  \omega_2 = \frac{I_1}{I_2}\,\omega_1 = 2\,\omega_1
  \qquad (\omega\text{ verdoppelt sich}).
\end{equation}

\textbf{Rotationsenergie:}
\begin{equation}
  \frac{E_\mathrm{rot,2}}{E_\mathrm{rot,1}}
  = \frac{\frac12 I_2\omega_2^2}{\frac12 I_1\omega_1^2}
  = \frac{(\frac12 I_1)(2\omega_1)^2}{I_1\omega_1^2}
  = \frac{\frac12\cdot 4}{1} = 2.
\end{equation}
Die Rotationsenergie verdoppelt sich ($E_\mathrm{rot}\propto 1/I$ bei
konstantem $L$: $E_\mathrm{rot}=L^2/2I$).

\textbf{Herkunft der Energie:} Beim Einziehen der Arme leistet der Sportler
gegen die Zentrifugalkraft positive Arbeit (die Armmuskeln ziehen die Masse
nach innen), die als zusätzliche Rotationsenergie gespeichert wird; die Zunahme
ist also $=W_\mathrm{Arme}=\Delta E_\mathrm{rot}=\frac12 I_1\omega_1^2$.

\begin{ergebnis}
$L$: Faktor $1$ (unverändert); $\omega$: Faktor $2$; $E_\mathrm{rot}$: Faktor
$2$ (Verdopplung). Die zusätzliche Energie kommt aus der Arbeit der Armmuskeln,
die beim Anziehen der Arme gegen die Zentrifugalkraft geleistet wird.
\end{ergebnis}

\begin{kontrolle}
\textbf{Formel-Check:} $E_\mathrm{rot}=L^2/2I$: halbiert man $I$ (konstantes
$L$), so verdoppelt sich $E_\mathrm{rot}$. \checkmark
\textbf{Energieerhaltung} in der Gesamtbilanz: Die Zunahme von $\Delta
E_\mathrm{rot}$ muss von der muskulären Arbeit kommen - kein Widerspruch,
da das System nicht isoliert ist (Arme als "Muskelmaschine").
\end{kontrolle}

\begin{merksatz}
Ohne äusseres Drehmoment bleibt $L=I\omega$ konstant; verkleinert man $I$, so
wächst $\omega$ reziprok. Die Rotationsenergie $E_\mathrm{rot}=L^2/2I$ wird
dabei grösser, und diese Energie muss von aussen (Muskelarbeit) zugeführt
werden.
\end{merksatz}